\cotefile{85}
\scanpage{16}
% L01–02
donc $(\mathbb D_\infty\times\mathbb Z_2)/N\simeq\mathbb D_\infty/N\cap\mathbb D_\infty$ et on applique encore le $=$ corollaire.

% L03–10
\uline{NB} Les groupes $\deleted{\mathbb D_{\unsure{i|\infty}},}\ \mathbb D_i\times\mathbb Z_2$ ($i$ infini, ou fini pair) et $\mathbb D_i$ sont-ils effectivement plongeables \unsure{sur|dans} un $H=N(T)$ ou $B$ ? \inserted{$k$ fixé (alg. clos de \unsure{car})}
% slanted
\margin{explicitons le cas des représentations \unsure{irr. diédrales|irr. diédriques|non diédrales}}
\deleted{Pour les $\mathbb D_i$, c'est évident que oui pour $N(T)$, \unsure{prouver que ce sont alg. clos|puisque ce sont alg. clos|pourvu que le seul alg. clos} les \unsure{ss groupes|gp. des}. le cas pour $\mathbb D_i\times\mathbb Z_2$ car déjà \unsure{cas} groupe de $\mathbb Z/2\mathbb Z\times\mathbb Z/2\mathbb Z$ vient pas} \inserted{\unsure{V$_4$|$\simeq V_4$}}

% L11–13
\deleted{\uline{Proposition} \inserted{Pour $k$ alg. clos, on envisage la \unsure{c.}} Pour que $\mathbb D_i$ \unsure{admette|admettent} une repr. \unsure{projective} dans $\mathrm{GP}(1,k)$, il faut que}

% L14–16
\uline{Proposition} Soit $k$ corps alg. clos (\deleted{de} car $p$), $i\in\mathbb N^{\deleted{\unsure{*}}}$, considérons la condition
\[
(\ast)\quad\exists\ \text{repr. fidèle}\quad\mathbb D_i\lhook\joinrel\longrightarrow\mathrm{GP}(1,k)
\]

% L17–20
\uline{Alors} :

a) Si $p=0$, \inserted{ou $i\deleted{\unsure{\geq 1}}=1$}, $(\ast)$ est satisfaite. \deleted{\unsure{Si $p\ne0$ et $i=1$, $(\ast)$ est satisfaite}.}

b) \deleted{Si $p\ne0$, (i.e. $\mathbb D_i=\mathbb D_0=\mathbb D_\infty$)} \inserted{Supposons $p>0$. Si $i=0$ (i.e. $\mathbb D_i=\mathbb D_0=\mathbb D_\infty$)} $(\ast)$ équivaut à la condition : \deleted{\unsure{$k$ n'est le c.}, \illegible{1}} $k$ non algébrique sur le corps premier (ou encore : $k^*$ contient un élément d'ordre infini)

% L21–25
\deleted{\unsure{b)|c)} Si $p\ne0$, et}

c) Si $i\geq2$, $(\ast)$ équivaut à la condition :

\deleted{On a}

$(\ast a)$ On a \deleted{$i$ premier à} $p\nmid i$, ou bien $i=p$, \deleted{\unsure{et le groupe} \illegible{1}}

\deleted{ou bien $i=2p$ et $p$ impair.}

% L26–32
Dém. a), b) immédiats, en \deleted{\unsure{$(\ast)$|cas}} que $i=0$ \unsure{$(\ast)$|car}
% slanted
\margin{\unsure{explicitons pour|explicitons un peu|explicitez ceux}}
équivaut à l'existence d'un él. d'ordre infini dans $k^*$ \{ou dans $k$.\}

\deleted{Pour $i\geq2$, on note que les}

\deleted{c) On sait que les sous-groupes \inserted{d'ordre fini} de $\mathrm{GP}(1)$ sont \unsure{ou les} les groupes \unsure{cycliques et premiers à $p$, ou les}}

\deleted{b) Si \unsure{$H(k)$} \unsure{$\simeq k$ ou}, \unsure{prenant} $H=\unsure{\mathrm V_{1/2}\,k|\mathrm V_{1,2}\,k|\gamma\underset{1/2}{\cdot}U}$ $(H=\unsure{\mathrm G_m}\text{ ou }\mathrm G_a)$}
% archive sheet number: 15
\scanpage{17}
% L01–03
Cas c)

\deleted{d'où $H(k)\simeq\gamma(k)\,\unsure{\ltimes|\cdot}\,U(k)$ $(U(k)\simeq k\text{ ou }k^*)$ et}

% L04–09
$(\ast)\Longrightarrow(\ast a)$ \deleted{si $p$ impair}, car les ordres des élts d'ordre fini de $\mathrm{GP}(1)$ sont les $i$ premiers à $p$, \deleted{et $p$ \illegible{1}} (correspondant aux él. $\frac12$ simples) et $p$ (correspondant aux éléments unipotents). \deleted{Mais \inserted{et} si $i=p=2$, \unsure{$G$} $\mathbb D_i=\mathbb Z_2\times\mathbb Z_2$ n'est pas contenu dans $\mathrm{GP}(1)$ (car \unsure{les uni} les $G_a$ \inserted{$\subset\mathrm{GP}(1)$} sont les unipotents maximaux). Et}

% L10–12
$(\ast a)\Longrightarrow(\ast)$ immédiat, en distinguant et examinant (pour le cas $i=p$) le cas $p$ impair, et le cas $p$ pair (alors cas chevauchant, l'image de $\mathbb D_2$ est à choisir dans $U$ !).

% L13–19
\deleted{\uline{Cor.} \unsure{Soit $p$, $i$ préc.|Soit $i$ pair} \unsure{Supposons $i$ pair.} Considérons}
\[
\deleted{(\ast')\quad\exists\ \text{repr. fidèle}\quad\mathbb D_i\times\mathbb Z_{/2}\longrightarrow\mathrm{GP}(1,k).}
\]
\deleted{a) Si $p=0$, \deleted{ou $i=1$} $(\ast')$ est satisfaite sss $i=0$}

\deleted{b) Si $i=0$, \inserted{donc $\mathbb D_i\times\mathbb Z_{/2}=\deleted{\unsure{\mathbb D_i}}=\mathbb D_\infty\times\mathbb Z_{/2}$}, $(\ast')$ est satisfaite sss \deleted{\unsure{$p$ est pair}} $p\ne2$.}

\deleted{c) Si $i\geq2$, $p>0$, alors $(\ast')$ n'est pas satisfaite.}

\deleted{\uline{Dém} \unsure{a) immédiat} (crit. de parité}

% L20–26
On s'intéresse aux repr. de $\mathbb D_\infty$ \inserted{$=\unsure{\mathrm G_{\Delta_1}|G_{P1}|G_1}=\{\underline\tau_0,\underline\tau_1\mid\underline\tau_0^2=\underline\tau_1^2=1\}$}, dans des $G$ \inserted{$\subset\uline{\mathrm{Aut}}(X)$} formes de $\mathrm{GP}(1)$ (sur $S$) \deleted{$\Longleftrightarrow$} la cat de celles-ci est équivalente à celle des représentations \deleted{\unsure{en}} projectives de dim $2$ de $\mathrm G_1$, \inserted{(dans un fibré projectif $X$)} \uline{munies de} \deleted{repr.} quelques bons invariants.
% slanted
\margin{On pose $\tau_i=\varphi(\underline\tau_i)$ $(i=0,1)$ : $u=\tau_0\tau_1=\varphi(\underline u)$}

% L27–31
\deleted{On les} \deleted{On va regarder} on va regarder d'abord le cas où la condition suivante est satisfaite

\deleted{$(\ast)$ $\forall s\in S$, $u_s$ $(=\tau_{0s}\tau_{1s})$ n'est}

\[
\begin{aligned}
(\ast)\quad &\forall s\in S,\qquad u_s^2\ne1\qquad
(\text{i.e. }\varphi_s:\mathbb D_\infty\longrightarrow G_s\ \text{ne se factorise}\\
&\hspace{18em}\text{pas par }\mathbb D_2)\\
&\text{i.e. }\tau_{0s},\tau_{1s}\text{ ne \uline{commutent pas}}
\end{aligned}
\]

% L32–35
\deleted{La catégorie de ces \unsure{fon} couples $(G,\rho)$ s'identifie à}

Or \deleted{\unsure{on sait}}, si on introduit $a_0,H_0,a_1,H_1$ comme \deleted{à la cat. des polygônes réguliers} d'habitude, la condition $(\ast)$ implique
\[
a_{0s}\notin H_1\quad(\Longleftrightarrow{}a_1\notin H_0)
\]
\inserted{\unsure{pour : le G inv.}} en tte fibre
% archive sheet number: 16
\scanpage{18}
% L01–06
donc loc. \inserted{fppf} ceci provient d'une situation d'un \uline{polygône régulier} \deleted{(projectif)} dans le plan projectif $X/S$, tel que $\deleted{\unsure{\beta}}\ \boxed{\beta\inserted{=\alpha+2}\ \text{inv.}}$ (ceci exprime $a_1\notin H_0$), \deleted{et} \inserted{ou encore tel} que la quadrique \deleted{\unsure{sur}} associée \inserted{$Q$} dans $\unsure{P|X}$ (l'unique quadrique) soit lisse. \deleted{\unsure{car}} Donc

\deleted{\unsure{la réf.} \illegible{1} $\unsure{\frac14|\frac12}\det\left(\begin{smallmatrix}0&&\\&2&\\&&2\end{smallmatrix}\right)$}

% L07–12
\uline{Th} \inserted{$S$ schéma.} a) Soit $C$ la catégorie des couples $(G,\varphi)$ d'une forme $G$ de $\mathrm{GP}(1)_S$ et d'un hom $\varphi:\mathbb D_\infty\longrightarrow G(S)$, \inserted{(d'où $\tau_0,\tau_1,u\in G(S)$)} satisfaisant à
\[
(\ast)\quad\forall s\in S,\qquad \deleted{\unsure{\{}}u_s^2\deleted{\unsure{\}}}\ne1 \qquad
(\text{i.e. }\varphi_s\text{ ne se factorise pas par }\mathbb D_2)
\]
% slanted
\margin{Dire qu'il existe un \uline{unique} div de degré $2$ fixe sous $\mathbb D_\infty$ -- qui correspond à un \unsure{sous-groupe} ``spécial'' de $G$, \unsure{lequel est ce|lequel est un} sous ? \unsure{Cf p.|t.p.}\ldots}

% L13–17
b) $\forall\ \beta\in\Gamma(S,\mathcal O_S^*)$, \deleted{\unsure{soit}} \inserted{soit} $U(\beta)$ le groupe \deleted{\unsure{forme}} lisse \deleted{Soit $C'$ la catégorie des couples formés d'un} de dim $1$ sur $S$, \deleted{défini} formé des $\lambda$ tels que $1+\lambda\beta$ inversible, avec la multiplication
\[
[\lambda]+_{\boldsymbol{.}}[\lambda']=[\lambda+\lambda'+\lambda\lambda'\beta]
\]
Soit $C'$ la catégorie des couples $(\alpha,R)$, où $\alpha\in\Gamma(S,\mathcal O_S)$ \uline{telle que $\beta:\alpha+2$} inv., et $R$ un torseur sous $\deleted{\unsure{G_{2-\alpha}}}\ U(2-\alpha)$.

% L18–23
c) Soit $\alpha\in\Gamma(S,\mathcal O_S)$, \deleted{\unsure{section de} $\mathcal O_S$} \inserted{et considérons} le polygône régulier principal de inv. $\alpha$, d'où, \inserted{si $\alpha+2$ inv. donc $Q_\alpha$ lisse}, une représentation
\[
\varphi_\alpha:\mathbb D_\infty\longrightarrow G_\alpha(S)
\]
\deleted{\unsure{$\ne$} telle \unsure{que}} \inserted{où $G_\alpha=\mathrm{Aut}(X_\alpha,Q_\alpha)$ et} \deleted{$U_{\unsure{\alpha}}$ opère sur $\mathrm{Aut}(G_\alpha,\varphi_\alpha)$.}
\[
U_{2-\alpha}\simeq\uline{\mathrm{Zentr}}(\varphi_\alpha)=\uline{\mathrm{Aut}}(G_\alpha,\varphi_\alpha)
\]
Si $R$ est un torseur sous $G_\alpha$, on peut l'utiliser pour former $(G_\alpha^R,\varphi_\alpha^R)$.

% L24–26
d) On trouve ainsi un foncteur $C'\longrightarrow C^{\unsure{*}}$, et \deleted{\unsure{X}}celui-ci est une équivalence de cat. Le foncteur quasi-inverse est
% archive sheet number: 17
\scanpage{19}
% L01–07
le foncteur $C\longrightarrow C'$ défini ainsi : donné $(G,\varphi)\in\mathrm{Ob}\,C$, on définit $\alpha$ comme l'invariant de $\varphi(\underline u)=u$, et $R$ comme le \deleted{fermé des} \unsure{sous foncteur|des foncteurs|faisceau} de $X$ (\inserted{plan projectif} enveloppe de la droite proj. $X/S$ définie par $G$) \deleted{fermé des} $H_1\setminus\{c\cup d\}$, qui est un torseur sous $U_{2-\alpha}\underset{\text{iso can}}{\sim}\deleted{\unsure{\mathrm{Aut}(G,\varphi)}}=\uline{\mathrm{Aut}}(X,\varphi)$

% L08–12
\ldots\ \unsure{Quant à|tenant à|tenant} : $\uline{\mathrm{Aut}}(\deleted{\unsure{Q}}X,\varphi)$, c'est le ss groupe \deleted{\uline{Corollaire} Soit en la \unsure{situation|relation}} de $U_{\beta\unsure{:=|=}2-\alpha}$ formé des $[\lambda]$ tels que
\[
\lambda(2+\lambda\beta)=0\quad\text{i.e.}\quad\deleted{\unsure{2\ne}}\ \lambda\mathbin{\ast_\beta}\lambda=1 .
\]
Donc

% L13–20, crossed by long diagonals
\deleted{1) Cas \unsure{$2\beta\ne0$|$2\rho\ne0$} \inserted{\unsure{$\beta=2-\alpha\ne0$|inversible}} \deleted{i.e. $\alpha\ne2$}, \deleted{i.e.} \inserted{i.e. $U_\beta\simeq$ un tore $\alpha\ne2)$ et $\beta\ne2$. \inserted{(car \unsure{inv.})} Alors} \deleted{\unsure{$U_\beta$ est tore. Alors}}}

\deleted{$\uline{\mathrm{Aut}}(X,\varphi)\simeq\deleted{\unsure{\mathbb Z_2}\ \unsure{\Leftrightarrow}\mu_2}\subset G_m$, \inserted{engendré par $\sigma$.} D'ailleurs l'image de $\deleted{\unsure{G_\alpha}}U_\alpha$ dans \deleted{Aut} $G\deleted{\unsure{(X)}}$ est l'unique tore maximal qui contient \deleted{\unsure{l'él.|l'image}} $u$, donc si $u$ est d'ordre fini pair $2n$, on a $\sigma=u^{\unsure{n|m}}$.}
% slanted
\margin{\deleted{\unsure{Dès que} \inserted{dans ces deux cas,} \unsure{dès qu'on pose|dès qu'on a} \unsure{on a} une représentation fidèle de $\mathbb D_{\unsure{\alpha|i|\infty}}\times\mathbb Z_2$ ($i$ pair) dont les \unsure{extensions|restrictions} à $\mathbb D_i^{\unsure{\alpha}}$ sont déter\unsure{.}}}

% L21–26
\deleted{2) Cas $\beta=0$, \deleted{$p\ne2$} \deleted{$2\ne0$} (i.e. $U_\beta\simeq G_a$ i.e. $\alpha=2$) \inserted{et \deleted{$p\ne2$} $2$ inv.} Alors, l'image de $U_\beta$ dans $G$ est l'unique \deleted{\unsure{un sous}} ssgroupe \deleted{unipotent} $\simeq G_a$ qui contient $u$. \deleted{\unsure{lorsque $p\ne$}} \deleted{\unsure{Si $2$ inv., toujours}} $\uline{\mathrm{Aut}}(X,\varphi)=\deleted{\unsure{\mu}}\ 1$}
% archive sheet number: 18
\scanpage{20}
% L01–04
\deleted{3) $\beta\ne0$, \inserted{inv.} $\deleted{\unsure{\beta|\rho}}p=2$ \inserted{et car $2$} (i.e. $U_\beta\simeq G_m$) Alors}

\[
\deleted{\mathrm{Aut}(X,\varphi)=\{1\}\qquad\uline{\mathrm{Aut}}(\unsure{X|G|Q},\varphi)\simeq{}_2U_\beta\overset{\text{can}}{\simeq}\mu_2}
\]
\deleted{On conclut à}

\uline{Proposition} $\uline{\mathrm{Aut}}(G,\varphi)\simeq {}_2(U_\beta)$

% L05–11
\uline{Corollaire 1} Supposons $\beta$ \inserted{$=2-\alpha$} inv, i.e. $U_\beta\overset{\text{loc}}{\simeq}G_m$, i.e.

\deleted{a) Si $\beta$ inv., et $2.1_S$ inversible, ou nul, alors}

\[
U_\beta\xrightarrow[\text{can}]{\sim}G_m,\qquad \inserted{\lambda\mapsto1+\lambda\beta}
\]
Alors : a) l'image $T$ de $U_\beta$ dans $G$ est l'unique tore maximal de $G$ \unsure{contnt|contenant} $u$. \deleted{Si \unsure{alors}}

b) $\uline{\mathrm{Aut}}(G,\varphi)={}_2T\overset{\text{can}}{\simeq}\mu_2$

\deleted{b) Si $2.1_S$ inv. ou $2.1_S=0$, Alors}

% L12–19
\uline{Corollaire 2} Supposons $\beta=2-\alpha=0$ i.e. $\alpha=2$, i.e. $U_\beta\overset{\text{loc}}{\simeq}G_a$ i.e. $\deleted{\unsure{U_\beta\ G.|U_\rho\simeq}}\ U_\beta\lhook\joinrel\longrightarrow\deleted{\unsure{G_m}}G_a$ \deleted{\unsure{$\lambda\mapsto\lambda$|$\lambda\mapsto1$}} est \unsure{multiplicatif|multiplicat.if}. Alors a) L'image $U$ de $U_\beta$ dans $G$ est l'unique \deleted{\unsure{ssgpe $\simeq G_a$}} ss groupe unipotent maximal de $G$ contnt $u$.

b) Si $2.1_S$ inv, Alors $\uline{\mathrm{Aut}}(G,\varphi)=\{1\}$

c) Si $2.1_S=0$ i.e. $S$ de car. $2$, Alors
\[
\uline{\mathrm{Aut}}(G,\varphi)=U\simeq G_a.
\]

% L20–24
\uline{Cor. 3} Soit $\gamma$ un gpe, et soit $\Gamma=\mathbb D_\infty\times\gamma$. Considérons la cat. $C(\gamma)$ des couples $(G,\varphi)$, $G$ forme de $\mathrm{GP}(1)_S$, $\varphi:\Gamma\longrightarrow G(S)$. Cette catégorie est équivalente à celle
% archive sheet number: 19
\scanpage{21}
% L01–04
des triples $(\alpha,R,\varphi_0)$, où a) $\alpha\in\Gamma(S,\mathcal O_S)$ tel que $\alpha+2$ inv (i.e. $\alpha(s)\ne-2\quad\forall s\in S$), b) $R$ un torseur sous \inserted{le groupe} $\deleted{\unsure{G}}\ U_{\beta=\alpha-2}$ c) $\varphi_0:\gamma\longrightarrow({}_2(U_\beta))(S)$ est un hom de groupes.

% L05–09
\uline{Proposition} Considérons \unsure{toujours|l'hom.|homomorphismes} \unsure{une|sur} \deleted{\unsure{$i$ entier pair,}} \deleted{\unsure{et soit}}

\deleted{\uline{Corollaire} Considérons $\Gamma=\mathbb D_i\times\mathbb Z_{/2}$ (\unsure{$i$ entier $\in(2\mathbb N)$}) et un corps alg. clos $k$, enfin la condition}

\[
\deleted{(\ast)\quad\exists\ \text{hom injectif}\quad\Gamma=\mathbb D_i\times\mathbb Z_{/2}\lhook\joinrel\longrightarrow\mathrm{GP}(1,k).}
\]
\deleted{a.) Si $i=0$, Alors $(\ast)$ est satisfaite sss}

% L10–15
\unsure{situation|solution} $\varphi:\mathbb D_\infty\longrightarrow G$ ($G$ forme de $\mathrm{GP}(1)$). \deleted{\unsure{Cas 1|Cond.}} Soit $n\in\mathbb N^{\deleted{\unsure{*}}*}$, \inserted{$n\geq3$}, considérons $\mathbb D_{n}\overset{\text{can}}{=}\mathbb D_i\simeq\mathbb D_\infty/\{\underline u^n\}$. Considérons les conditions (équivalentes)

a) $u^n=1$, et $\forall s\in S$, $u_s$ d'ordre \unsure{$n$} exact \deleted{[Si $n=0$, il faut dire : $u_s$ d'ordre $\infty$]}

b) $\varphi$ se factorise par un \uline{mono} $\mathbb D_{i_S}\lhook\joinrel\longrightarrow G$.

% L16–22
On veut les exprimer en termes de l'invariant $\alpha$. On trouve : \deleted{\unsure{on est dans les cas}}

\deleted{1°) Si $n=\infty$ : il f. et s. que $\forall s\in S$, on ait \unsure{que a) et c)} $\forall s\in S$ tel \unsure{1°) où}}

c) 1°) $n.1_S$ est inversible, \uline{ou} $n$ \inserted{est} \deleted{\unsure{$=2p$ où $p$|$=p$} \unsure{premier}} \inserted{un} \deleted{ou bien $n$ \unsure{d'un $p$ premier dans} $\Gamma(S,\mathcal O_S)$} est premier impair \deleted{\unsure{inversible}}, \deleted{\unsure{alr}} \uline{et} 2) $\Psi_n(\alpha)=0$

[$\Psi_n\in\mathbb Z[U]$ polynôme $\frac12$ cyclotomique d'indice $n$]
% archive sheet number: 20
\scanpage{22}
% L01–04
\uline{Cor} Soit $k$ corps alg. clos, \inserted{$i$ entier \uline{pair} $\geq0$,} Considérons la condition
\[
(\ast)\quad\exists\ \text{hom injectif}\quad\Gamma=\mathbb D_i\times\mathbb Z_{/2}\lhook\joinrel\longrightarrow\mathrm{GP}(1,k)
\]
\uline{Alors}

% L05–11
a) Si $i\unsure{>|{\geq}|\ne|=}0$, \inserted{donc $i\geq2$}, la condition \deleted{$(\ast)$ n'est pas} \inserted{$(\ast)$ est} vérifiée sss \unsure{car.} $k=2$.
% slanted
\margin{distinguer le cas $i=2$ dans dém.}

b) Si $i=0$, elle est vérifiée sss ou bien car $k=0$, ou \deleted{bien} car. $k>0$ \uline{et} $k$ non alg. sur le corps premier.

\deleted{a) Si $i=2$, \inserted{donc $\Gamma\simeq(\mathbb Z/2)^3$}, la condition $(\ast)$ est vérifiée sss car $k=2$.}

% L12–18
a) Supposons $i\ne0,2$ i.e. $i\inserted{=2i'}>4$. La condition pour qu'il existe mono $\mathbb D_i\mathbin{\unsure{\subset|-|\hookrightarrow}}\mathrm{GP}(1,k)$ est que $p\nmid i$ ou $\deleted{\unsure{p=i|p\mid}}\ i=p$ $(\ne2)$, dans le premier cas $\alpha\ne2$ i.e. $\beta\ne0$, donc $\mathrm{Aut}(G,\varphi\mid\mathbb D_i)\simeq\mu_2$ \deleted{\unsure{$=\mu$}} donc l'image du gén. \inserted{$\sigma$} du facteur $\mathbb Z_{/2}$ de $\Gamma$ est soit $1$, \deleted{\unsure{car|on}} soit $u^{\unsure{i'|n|\nu}}$. d'où le résultat.

% L19–22
Dans le 2ième cas $\alpha=2$, i.e. $\beta=0$, et en car. $\ne2$ on a $\mathrm{Aut}(G,\varphi\mid\mathbb D_i)=(1)$, (donc un hom $\mathbb D_i\times\mathbb Z_{/2}$ s'annule sur $\sigma$) et en car. $2$\qquad\rule{3cm}{0.3pt}
% archive sheet number: 21
\scanpage{23}
% L01–02, title at upper right
\margin{(9 p)}
\uline{Invariants} des sections des formes de $\mathrm{GP}(1)$ etc
% archive sheet number: 22
\scanpage{24}
% L01 heading
Invariant\unsure{s} \deleted{\unsure{des}}d'un automorphismes d\deleted{e}'une droite projective.

% L02–06
a) Soit $G$ une forme de $\mathrm{GP}(1)$, ou de $\mathrm{Sl}(2)$ \inserted{sur $S$}, -- i.e. un schéma en groupes $\frac12$ simple de rang $1$ sur schéma $S$. Considérons le foncteur
\[
\mathbb I(G)=\underline{\mathrm{Hom}}_{\rm sch}(G,\mathbb E^1_S)^{\unsure{\leftrightarrows|\leftrightarrow|\circlearrowleft}}
\overset{\text{déf}}{=}\underline{\mathrm{Hom}}_{\rm sch}(G,\mathbb E^1_S)^G
\]
où $G$ opère sur le $\underline{\mathrm{Hom}}$ \deleted{via} autom. int. sur $G$.

% L07–14
Ce foncteur est représentable, \inserted{(vrai pour tt $G$ $\frac12$ simple)} et on a un isomorphisme canonique
\[
\mathbb I(G)\xrightarrow{A}\mathbb E^1_S
\]
qui, en tant que morphisme, peut être caractérisé par la propriété suivante : Si $G_m\overset{i}{\lhook\joinrel\longrightarrow}G$ est un monomorphisme, alors
\[
A(i(\lambda))=\lambda+\lambda^{-1}
\]
et itou après tout chgt de base.

% L15–21
\deleted{\uline{Proposition} Soient $u,v$ deux sections de $G$, telles que \inserted{on suppose} \inserted{\unsure{a) Donc que $u$|2) Prouvons que}}}
\deleted{\unsure{ou}} On a donc
\[
A(1)=2.
\]
[ \deleted{Supposons $G$ forme de $\mathrm{GP}(1)$, et} \inserted{Soit} \deleted{et si} $\sigma$ \deleted{est} une section de $G$ sur $S$ telle que $\sigma^2=\mathrm{id}$, $\sigma_s\ne\mathrm{id}_s\quad\forall s\in S$, \deleted{$\sigma$ soit d'ordre $2$ \unsure{exact} (donc ce qu'on peut appeler une réflexion de $G$)} alors
\[
A(\sigma)=-2
\]
% slanted
\margin{\uline{NB} Dans le cas d'une forme de $\mathrm{GP}(1)$, $\sigma$ est une ``réflexion'', dans le cas de $\mathrm{Sl}(2)$, en car rés. $\ne2$, \unsure{c'est la} section génératrice du centre\ldots}

% L22–26
C'est clair si $2$ inv. sur $S$, car alors $\sigma$ est contenu dans un \deleted{(unique)} tore maximal, \inserted{(unique dans le cas de $\mathrm{GP}(1)$ -- dans le cas de $\mathrm{Sl}(2)$, $\sigma$ est central)} qu'on peut supposer scindé (par localisation), et on a alors $\lambda=-1$, d'où $A(u)=-1+1/{-1}=-2$. Pour le cas général, on fait un argument de coeff universels \ldots. ]

% L27–28
\deleted{\uline{Prop} Dire que} Supposons que $u$ soit \uline{unipotent}, par quoi on entend qu'il existe loc. fpqc une \deleted{forme}
% archive sheet number: 23
\scanpage{25}
% L01–05
monomorphisme
\[
\mathrm{Aff}(1)_S=G_m\underset{1/2}{\boldsymbol{.}}G_a\xrightarrow{i}G
\]
et une section $t$ de $G_a$, telle que $u=i(t)$. Je dis qu'alors aussi on a
\[
A(u)=2
\]

% L06–12
En effet, on trouve plus \deleted{précisément} \inserted{généralement}
\[
A(i(\lambda:t))=\lambda+\lambda^{-1}
\]
On le voit p. ex. en notant que (pour $\lambda,t$ indét)
\[
f(\lambda,t)\overset{\text{déf}}{=}A(i(\lambda,t)),
\]
doit être un $\in\underline O_S[\lambda,t][\lambda^{-1}]$ satisfaisant les conditions
\[
\begin{aligned}
&\left\{\begin{aligned}
&f(\lambda,0)=\lambda+\lambda^{-1}\\
&\text{plus gén}\quad f(\underbrace{(\unsure{0|\omega},t)(\lambda,0)(0,-t)}_{(\lambda,0)(\lambda^{-1}t-t)=(\lambda,\lambda^{-1}t-t)})=\lambda+\lambda^{-1}
\end{aligned}\right.
\end{aligned}
\]

% L13–17
donc, posant $s=\lambda^{-1}t-t$ i.e. $t=\dfrac{s}{\lambda^{-1}-1}$, on a
\[
f(\lambda,s)=\lambda+\lambda^{-1}
\]
dès que $\lambda^{-1}-1$ inversible -- donc par prolongement, pour tt système $(\lambda,s)$.

% L18–24
Une condition \uline{nécessaire} pour que $u$ soit unipotent est donc que $A(u)=1$. On voit aisément que cette condition est suffisante au dessus d'un corps (utilisant le fait que sur un corps alg. clos, tt élément de $G$ est contenu dans un ss-groupe de Borel). J'ignore ce qui en est dans le cas général.
% archive sheet number: 24
\scanpage{26}
% L01–06
\uline{Proposition} Soient $u,u'$ deux sections de $G$, \deleted{telles} \inserted{on} suppose $A(u)-2$ \inserted{et $A(u)+2$} inv. [i.e. $u_s^2$ non unipotent $\forall s\in S$]. Pour que $u$ et \unsure{$v$|$u'$} soient conjugués loc. fpqc (ou encore, conjugués si $G$ forme de $\mathrm{GP}(1)$) il f. et \uuline{s.} que $A(u)=A(u')$.
\margin{\unsure{Ceci|Celà} reste vrai sans restriction sur \unsure{$\alpha(u)$|$A(u)$}, si $S$ spectre d'un corps et $u_s$ non central \unsure{$K_s$|$\forall s$})}

% L07–12
\deleted{En effet,} \inserted{Néc. évidente, prouvons suff.} L'hypothèse \inserted{sur $u$ et $u'$} implique que $u,v$ sont contenus chacun dans un (unique) tore maximal. Utilisant le th. de conjugaison, on peut supposer que c'est le $=$, et par localisation, que c'est $G_m$. Donc $u,u'$ correspondent à des sections $\lambda,\lambda'$ de $G_m$, et l'hypo $A(u)=A(u')$ signifie

% L13–18
\[
\lambda+\lambda^{-1}=\lambda'+\lambda'^{-1}\qquad\text{soit }\alpha\in\Gamma(S,\unsure{\underline O_S|\mathcal O_S})\tag*{$(*)$}
\]
ce qui signifie que $\lambda$ et $\lambda'$ sont \deleted{\unsure{racines}} \inserted{solutions} de la $=$ équation
\[
\deleted{\unsure{x}}t^2-\alpha t+1=0\tag*{$(**)$}
\]
dont le discriminant est
\[
\alpha^2-4=(\underbrace{\alpha-2}_{\text{inv.}})(\alpha+2)
\]

% L19–28
\deleted{il n'est donc nul qu'aux pts où $\alpha(s)=-2$.} \inserted{donc inv. par hypo} \deleted{Si $\alpha^2-4$ inv. i.e. $\alpha+2$ inv. i.e. $\alpha(s)\ne-2$ en tt $s\in S$, alors} Cette équation définit un rev. \uline{étale} de $S$, et \deleted{on trouve que l'on a loc deux solutions distinctes} deux sections sont soit égales loc, soit conjuguées par l'involution canonique $t\longmapsto t^{-1}$, on gagne donc. Il reste à faire une étude locale \deleted{\unsure{aux}} plus serrée aux pts où $\alpha(s)=\pm2$, i.e. $\lambda_s=\lambda'_s=\pm1$, prenons le cas $-1$ p.ex.
% archive sheet number: 25
\scanpage{27}
% L01–06
Les réductions habituelles nous ramènent au cas d'un anneau \inserted{de base} \deleted{artinien $\mathbb Z/p^n\mathbb Z$ $A$, corps} artinien $A$, corps résiduel \deleted{\unsure{$k=$|$k$}} $\mathbb Z/p\mathbb Z$, à prouver que si $\lambda\deleted{\unsure{\ne|=}}{,}\lambda'\in A^*$ sont tels que $\lambda+\lambda^{-1}=\lambda'+\lambda'^{-1}$, alors on a $\lambda'=\lambda^{\deleted{\unsure{\varepsilon}}}$ ou $\lambda'=\lambda^{-1}$. Supposons que l'on sache $\lambda'_n=\lambda_n$ \unsure{prouvons|prouver}
\[
\underset{\substack{\Vert\\\lambda'}}{\lambda'_{n+1}}=\underset{\substack{\Vert\\\lambda}}{\lambda_{n+1}}^{\pm1}\ ?
\]

% L07–11
On a \deleted{$\unsure{\lambda'_n=\lambda_n}$} $\lambda'=\lambda+\mu$, avec $\unsure{\mathfrak m|m}\mu=0$, \deleted{d'où} donc
\[
\begin{aligned}
\lambda'&=\lambda(1+\lambda^{-1}\mu)=\lambda(1-\mu)=\lambda+\mu\\
\lambda'^{-1}&=\lambda^{-1}\underbrace{(1-\mu)^{-1}}_{1+\mu}=\lambda^{-1}(1+\mu)=\lambda^{-1}\unsure{-|+}\mu\\
\lambda'+\lambda'^{-1}&=\lambda+\lambda^{-1}+\mu(\lambda^{-1}-\lambda^{\deleted{\unsure{-1}}})=\lambda+\lambda^{-1}.
\end{aligned}
\]
% slanted
\margin{\uline{NB} Il s'agit apparemment d'une propriété générale des morphismes \deleted{\unsure{qui}} ``ramifiés'' (non nets) en un point : \unsure{certaine|certains|existence} ``\unsure{continu|continuum}'' de \unsure{l'ens.} des solutions infinitésimales (dès le $1^{\rm er}$ ordre\ldots)}

% L12–16
Donc l'hypo $\lambda'+\lambda'^{-1}=\lambda^{\deleted{\unsure{-1}}}+\lambda^{-1}$ résulte de $\lambda'_n=\lambda_n$ i.e. $\lambda'=\lambda+\mu$, $\mathfrak m\mu=0$ \unsure{$\}$|?} ; a-t-on pour autant
\[
\lambda'=\lambda\quad\text{ou}\quad\lambda'^{-1}=\lambda\quad\deleted{\unsure{?}}?
\]
\deleted{i.e} \unsure{Or|On} ces relations signifient, \deleted{soit} \inserted{l'une} $\mu=0$,

% L17–23
\deleted{soit} \inserted{l'autre} \deleted{$\unsure{\lambda'=\lambda(1+\mu)=\lambda|\lambda^{-1}(1+\mu)=\lambda}\quad\text{i.e.}\quad\lambda\mu=0$}
\[
\begin{aligned}
&\deleted{\lambda^{-1}(1+\mu)}\quad\deleted{\lambda-\lambda^{-1}+\mu=0}\\
&\mu=\lambda^{-1}-\lambda
\end{aligned}
\]
donc ne correspondent qu'à deux valeurs possibles \inserted{au plus} de $\mu$ -- donc (si $k$ n'est pas réduit à $2$ él) on peut choisir $\mu$ de façon à mettre en défaut la relation $[(\lambda'=\lambda)\text{ ou }(\lambda'=\lambda^{-1})]$.
% archive sheet number: 26
\scanpage{28}
% L01–05
\uline{Remarque} Soit $j:G_m\mathbin{-}G$ un hom de groupes, pour néc. \deleted{injectif} \inserted{nous} supposons que le noyau soit plat, donc $j$ se factorise en
\diagram{\begin{tikzcd}[ampersand replacement=\&]
G_m \arrow[r,dash,"\lambda\to\lambda^n"] \arrow[rr,bend right=16,"j"'] \& G_m \arrow[r,dash,"i"] \& G
\end{tikzcd}}
On en conclut que
\[
A(j(\lambda))=\lambda^n+\lambda^{-n}=G_n(\lambda+\lambda^{-1})
\]

% L06–10
où
\[
G_n\in\mathbb Z[U]
\]
est un polynôme unitaire universel de degré $n$, [celui exprimant $T^n+T^{-n}\in\underset{\substack{\unsure{\wr|\simeq}\\\mathbb Z[U]}}{\mathbb Z[T,T^{-1}]^\sigma}$ $(\sigma T=T^{-1})$ \deleted{\unsure{comme}} $U=T+T^{-1}$]

% L11–15
% slanted
\margin{\uline{\unsure{remords|remonte|remarque}}}
\uline{Exemples} 1) Soit
\diagram{\begin{tikzcd}[ampersand replacement=\&]
G_m \arrow[r,hook,"i"] \& \underset{\substack{\text{forme de}\\\mathrm{Sl}(\unsure{1|2})}}{\widetilde G} \arrow[r,"p"] \& G=\widetilde G/\underset{\substack{\text{centre de}\\\widetilde G}}{\mu_2}
\end{tikzcd}}
on a donc
\[
\begin{aligned}
A(i(\lambda))&=\lambda+\lambda^{-1}=\alpha\\
A(pi(\lambda))&=\lambda^2+\lambda^{-2}=\alpha^2-2
\end{aligned}
\]

% L16–20
Plus généralement, indépendamment de \unsure{$i$|$A$} \unsure{il} \ldots{} pour \deleted{2) Soit $u$ une ps. réflexion dans $V$ (\unsure{fibré vctriel} de rg $2$ sur $S$)} toute section $u$ de $\widetilde G$ sur $S$
\[
\left\{\begin{aligned}
A'(p(u))&=A(u')=A(u)^2-2\\
A(u)&=\mathrm{Tr}\,u
\end{aligned}\right.
\]
[se ramener au cas précédent]

% L21–27
2) Soit $V$ Mod. loc. libre rg $2$ sur $S$, \inserted{$C=\unsure{\check{\mathbb P}|\breve{\mathbb P}}(C)$,} $u$ autom. de $V$, \unsure{:|;|,} \deleted{alors} (pas néc dans $\mathrm{Sl}(V)$) $u_C$ l'autom. de $C$ qu'il définit, alors on a
\[
A(u^C)=\frac{(\mathrm{Tr}\,u)^2}{\det u}-2
\]
(se ramener au cas 1), en écrivant loc $u=tv$, $t$ unité, $v\in\mathrm{Sl}(V)$\ldots)

% L28–29
Si p.ex. l'équation aux val. propres de $u$
\[
t^2-(\mathrm{Tr}\,u)t+\det u=0
\]
% archive sheet number: 27
\scanpage{29}
% L01–04
se décompose en $(t-\lambda_1)(t-\lambda_2)$, i.e.
\[
\lambda_1+\lambda_2=\mathrm{Tr}\,u,\qquad\lambda_1\lambda_2=\det u
\]
alors on trouve
\[
A(p(u))=\frac{\lambda_1}{\lambda_2}+\frac{\lambda_2}{\lambda_1}=\frac{\lambda_1^2+\lambda_2^2}{\lambda_1\lambda_2}
\]

% L05–08
P.ex. si $u$ est une \uline{pseudoréflexion}, \deleted{\unsure{de valeur}} de paramètre $\lambda=\det u$ -- on a les ``val. propres'' $\lambda$ et $1$, et on trouve
\[
A(pu)=\lambda+\lambda^{-1}\qquad(u\text{ ps. }\unsure{\text{réflex}|\text{réfl.}}\ \unsure{\text{pour}|\text{para}|\text{param}|\text{par}}\ \lambda)
\]

% L09–13
On peut aussi le voir en notant qu'une ps. réflexion s'insère dans une famille à un paramètre, correspondant à
\diagram{\begin{tikzcd}[ampersand replacement=\&]
G_m \arrow[r,hook] \arrow[rr,bend right=17,dash,"\text{mono}"'] \& \uline{\mathrm{Gl}}(V) \arrow[r] \& \uline{\mathrm{Aut}}\,\deleted{\unsure{(C)}}\inserted{C}
\end{tikzcd}}
et appliquer la définition générale de l'invariant.

% L14–18
\uline{Proposition} Soit $u$ une section de $G$, \deleted{Pour que} \inserted{telle que} $u_s\ne\mathrm{id}_s$, $\forall s\in S$. \deleted{$u^2=1$, il suf.} Si $u^2=1$, \inserted{et supposant car. rés. $\ne2$ si $G$ forme de $\mathrm{Sl}(2)$,} on a
\[
A(u)=-2,
\]
la réciproque étant vraie si $S$ est réduit.

% L19–26
\uline{Dém} Dans le cas \deleted{de} \inserted{$2$ inv., pour} une forme de $\mathrm{GP}(1)$, l'hypothèse \inserted{$u_s\ne1_s$ \deleted{\unsure{$\forall s\in S$}}, $u_s^2=\mathrm{id}\ \forall s\in S$} implique que $u$ appartient à un unique tore maximal -- pour une forme de \deleted{ss grpe canonique de} $\unsure{\mathrm{Sl}(2)|\mathrm{Sl}(1)}$, que c'est la section gén. du centre $\mu_2\simeq\{\pm1\}_S$, donc en tous cas loc $u$ \deleted{\unsure{est une}} \inserted{est} une section d'un tore maximal $T\simeq G_m\mathbin{\unsure{\subset|-|\hookrightarrow}}G$, correspondant à la section $-1$ de $G_m$, d'où $A(u)=-2$.

% L27
Dans le cas d'une forme de $\mathrm{GP}(1)$, on \unsure{ne|en}
% archive sheet number: 28
